\chapter{INTRODUCTION}

\label{ch:introduction}

\section{Rationale}

Planning a university degree is a constraint-satisfaction problem that students usually solve
with a printed curriculum sheet and a spreadsheet. The sheet states which courses exist, which
academic year and semester each is nominally offered in, and occasionally which courses must
precede which; the spreadsheet records what the student has already passed. Neither artifact can
answer the questions students actually ask: which courses am I now eligible for, what happens to
the rest of my plan if I drop this one, which elective group should claim this course, and do I
still qualify for the thesis track.

At the International University -- VNU HCMC the problem is sharper than the generic case for three
reasons. First, prerequisite relationships are declarative but numerous, so eligibility depends on
a graph rather than a list. Second, the curriculum repeats courses across elective groups: one
course can appear in several groups and in more than one semester, so a completed course must be
attributed to exactly one group without being counted twice. Third, the final year branches: a
student above a grade threshold takes the thesis, a student below it takes additional coursework
instead, which changes both the number of remaining courses and the shape of the last semesters.

A fourth reason emerged while building the system. Much of the information a planning tool needs is
\emph{per program}, not global. The same calculus course is a first-semester course for computer
science and a second-year course for business, and the prerequisite edges around a shared course
differ between majors. A single global catalogue therefore cannot represent more than one program
honestly, and a tool that quietly unions two programs would produce prerequisites that neither
program declares.

\section{Problem Statement}

The thesis addresses one question: how can a web-based system represent a university program, the
student's record, and the institution's resource constraints precisely enough that a student can
plan a semester and an administrator can reason about capacity, while making no claim that the
available evidence does not support?

The question decomposes into six design problems:

\begin{enumerate}
  \item \textbf{Program representation.} Model program membership, per-program placement and
        per-program prerequisites without duplicating shared courses or merging programs.
  \item \textbf{Account scope.} Bind every read and write to one account and one program context,
        so a stale client cannot write against the wrong scope.
  \item \textbf{Progress integrity.} Keep completion consistent: a completion is refused when its
        prerequisites are unmet, and removing a prerequisite must remove everything that depended
        on it, transitively, on the server.
  \item \textbf{Grade and difficulty evidence.} Record attempts immutably, compute a defensible
        grade average, and derive course difficulty from student votes in a way that stays honest
        when few votes exist.
  \item \textbf{Saved-plan integrity.} Let a student edit a saved semester without admitting an
        unsupported, duplicated or unprincipled course selection, and without losing work when a
        reply is lost.
  \item \textbf{Planning and capacity under uncertainty.} Offer semester recommendations and an
        administrative allocation simulation that state their assumptions, degrade explicitly when
        data is missing, and never present a simulated allocation as a registration decision.
\end{enumerate}

\section{Objectives}

The objectives of this work are:

\begin{enumerate}
  \item To design and implement a program-aware curriculum model in which courses are global but
        membership, placement and prerequisites are per program, with an imported reference
        curriculum that preserves repeated elective occurrences and uncoded requirements.
  \item To implement cookie-based accounts in which the server re-reads the current role and
        program context for every request, and in which every mutation carries an expected scope
        that is compared inside the same transaction as the write.
  \item To implement server-authoritative completion with mandatory prerequisite validation and a
        transactional transitive cascade on uncompletion.
  \item To record numeric grade attempts immutably, compute a credit-weighted average from the
        highest score per course with physical training excluded, and record per-semester averages
        from stored term metadata without inferring dates.
  \item To collect completion-gated course difficulty ratings, combine them with a per-program
        prior through Bayesian shrinkage, and display the vote count beside every estimate.
  \item To implement a saved-semester editor with atomic validation, immutable request receipts,
        recoverable unsent drafts, per-course blocker explanations and an explicit display order.
  \item To implement an administrator simulation of resource settings, cohort demand, capacity and
        two allocation models, with versioned contracts, private per-student outcomes and an
        explicit single-job executor rather than an unattended daemon.
  \item To verify the artifact with fail-first tests against a real PostgreSQL database, a single
        reproducible continuous-integration gate, and a measured accessibility pass in a real
        browser.
\end{enumerate}

\section{Scope and Limitations}

The artifact is a work in progress toward an institutionally valid planner, and the thesis is
explicit about where that line sits.

\textbf{Within scope.} A single reference program (Computer Science) imported from the program's
published curriculum page; a second tier of program support in the data model, exercised by tests
and by the administrator simulation, but not populated with verified Information Technology or
Data Science data. Simulated students, simulated resource configurations and simulated allocations.
Numeric grading on a 0--100 scale. Engineering evidence: automated tests, continuous integration,
reproducibility, measured profiling of the allocation models, and a browser accessibility audit.

\textbf{Outside scope.} Institutional validation of any program: the imported curriculum is a
reference, and neither its credit total nor its elective rules are asserted to be the signed
policy of a cohort. Course offerings, meeting times, rooms, laboratory suitability and staff
qualification, none of which were available, so calendar validity is not claimed. Editing of a
student's record inside an assigned program context, which requires context-keyed client caches
and a claim-transfer policy that were deliberately left out. A real deployment with a public
origin, a production database and operational recovery. A user study: no students were recruited,
and no claim is made that recommendations improve outcomes.

\textbf{Known open finding.} The server suite occasionally returns an authentication refusal where
a business-level refusal is expected, and occasionally fails a suite that passes in isolation. The
cause has not been established. It is recorded rather than suppressed, and no assertion or timeout
was weakened to obtain a passing run.

\section{Research Framework}

The work follows the Design Science Research Methodology of Hevner et al.\ \cite{hevner2004},
which suits an artifact whose contribution is evaluated by construction, demonstration and
measurement rather than by hypothesis testing on a population. The six activities are mapped in
Chapter~\ref{ch:methodology} to the increments actually delivered, and the evaluation is confined
to what the artifacts of those increments can support.

\section{Structure of Research}

Chapter~\ref{ch:literature} reviews degree-audit and planning systems, prerequisite-graph
analysis, course sequencing and timetabling, difficulty estimation, and web accessibility.
Chapter~\ref{ch:methodology} states the methodology, the requirements derived from it, the
architecture, the data model and the verification strategy.
Chapter~\ref{ch:implementation} describes the implementation increment by increment, including the
decisions that were deliberately reversed during development.
Chapter~\ref{ch:results} demonstrates the system and reports the verification evidence, the
measured performance of the simulation models, and the accessibility results, together with the
threats to validity.
Chapter~\ref{ch:conclusion} states the contributions, the limitations and the work that remains.
